\begin{textsample}{NEG Dim 6 – global\_south\_2024\_05 – Score 1.00 – global\_south\_2024\_05/108616729.txt}  \label{ex:f6_neg_294}
Xi pledges more Gaza aid and talks trade at summit with Arab leaders CHINESE President Xi Jinping, right, and Bahrain ' s King Hamad bin Isa Al Khalifa arrive to attend the opening ceremony of the 10th ministerial meeting of the China-Arab States Cooperation Forum at the Diaoyutai State Guesthouse in Beijing on Thursday, May 30, 2024.
BEIJING -- Chinese President Xi Jinping reiterated calls for the establishment of an independent Palestinian state and promised more humanitarian aid for people in Gaza as he opened a summit with leaders of Arab states Thursday in Beijing."Since last October, the Palestinian-Israeli conflict has escalated drastically, throwing people into tremendous suffering,"Xi said in a speech opening the China-Arab States Cooperation Forum."War should not continue indefinitely."He restated China ' s backing of a two-state solution and pledged 500 million yuan ( \$69 million ) in humanitarian aid for Gaza.
He also promised to donate \$3 million to a United Nations agency that provides assistance and relief to refugees of the Israel-Hamas war.
Beijing and the Arab states back @ @ @ @ @ @ @ @ @ @ international condemnation after its strike in the southern Gaza city of Rafah in which at least 45 people were killed over the weekend.
The overall Palestinian death toll in the war exceeds 36,000, according to the Gaza Health Ministry.
Beijing has long backed Palestinians and denounced Israel over its settlements in the occupied territories.
It has not criticized the initial Hamas attack on October 7 -- which killed about 1,200 people -- while the United States and others have called it an act of terrorism.
However, China has growing economic ties with Israel.
Besides addressing the war, Xi also called on Arab states to deepen cooperation in areas such as trade, clean energy, space exploration and health care.
The summit attended by heads of state from Egypt, the United Arab Emirates, Bahrain and Tunisia among others was set to focus on China ' s expanding trade ties and on security concerns related to the Israel-Hamas war."China ' s priorities in the region are primarily economic,"said Maria Papageorgiou, a @ @ @ @ @ @ @ @ @ @."It wants to continue the momentum established in recent years with Gulf States and expand its investments, particularly in trade, technology ( 5G networks ), and other cyber initiatives."Additionally, China wants to present itself as an alternative to the West and a more credible partner to the region, one that doesn't interfere in the nations ' domestic affairs nor exert pressure, Papageorgiou said.
Present at the forum is Egyptian President Abdel-Fattah el-Sissi, who met Xi on Wednesday.
The two leaders signed a series of cooperation agreements in areas such as infrastructure, technology and food imports meant to further their countries ' ties.
China has invested billions of dollars in Egyptian state projects, including a Suez Canal economic zone and a new administrative capital east of Cairo.
Investments between Egypt and China amounted to around \$14 billion in 2023, compared to \$16.6 billion in 2022, according to Egypt ' s statistics agency.
Also at the forum are Tunisia ' s President Kais Saied, Emirati President Sheikh @ @ @ @ @ @ @ @ @ @ The China-Arab States Cooperation Forum was established in 2004 as a formal dialogue mechanism between China and Arab states.
China is Tunisia ' s fourth-largest trading partner after Germany, Italy and France.
Beijing has financed hospitals and sports complexes in Tunisia, and its companies have been contracted to build strategic infrastructure such as bridges and deep-water Mediterranean ports.
The UAE also has expansive, growing economic ties with China and has faced US criticism for an alleged Chinese military facility being built in Abu Dhabi.
Besides China ' s expansive trade ties in the Middle East, it has increasingly sought to play a diplomatic role in the region.
In 2023, Beijing helped broker an agreement that saw Saudi Arabia and Iran reestablish ties after seven years of tension in a role previously reserved for longtime global heavyweights like the US and Russia.

% matched lemmas: 
\end{textsample}
